% Supplementary report section.  Requires amsmath, amssymb, amsthm, hyperref.
% Ordinary mathematical proof with an exhaustive finite certificate.
\section{A complete obstruction to propositional reconstruction of the fourth input}
\label{sec:all-fourth-input-obstruction}

The current AX constructor uses the four-parameter implicational schema
\[
 L(a,b,c,u)=((c\to u)\to a)\to((a\to c)\to(b\to c)).
\]
The fourth input $u$ requires a separate read from the proof stream.
The preceding incoming companion measured smaller source emitters obtained
by replacing it with an already available scalar or a short formula.
Its logical obstruction covered every positive implication term and one
particular negative term.  The following two-matrix argument covers
\emph{every} finite implication term in the remaining inputs, including
terms containing falsity.  This is an obstruction to recovering the
classical logical subsystem by this particular source simplification;
it is not a statement that the whole modified set theory has the wrong
consistency strength.

\subsection{The exact subsystem}

Use the named-variable formula language with membership, native equality,
implication, universal quantification, and falsity $F$.  As in the
published machine, abbreviate
\[
 \neg A=(A\to F),\qquad
 \exists x\,A=\neg\forall x\neg A,
 \qquad\top=(F\to F).
\]
Retain modus ponens, generalization, and the seven other logical schemas
from the
\href{https://github.com/jeff87654/busy-beaver/blob/zf-254/ZF/m254/design/logic.md}
{upstream logical specification}:
\begin{align*}
 \rho_1&:\quad \exists x\,(x=y),\\
 \rho_2&:\quad
  \forall x(A\to B)\to(\forall x A\to\forall x B),\\
 \rho_3&:\quad x=y\to(z\in x\to z\in y),\\
 \rho_4&:\quad \forall x\forall y A\to\forall y\forall x A,\\
 \rho_5&:\quad x=y\to(x\in z\to y\in z),\\
 \rho_6&:\quad \exists x\forall x A\to A,\\
 \rho_7&:\quad x=y\to(x=z\to y=z).
\end{align*}
Also retain both fresh atomic vacuity schemas
\[
 P\to\forall xP\quad(x\text{ does not occur in the membership or equality atom }P)
\]
and irreflexivity $\forall x\neg(x\in x)$.
For the theorem we may additionally give the proposed replacement the
full schemas
\[
 K:\ A\to(B\to A),\qquad A\to A,
 \qquad F\to A,\qquad\neg\neg A\to A.
\]
Adding these makes the nonderivability statement stronger.
Neither the container axiom CP nor any instance of GIC is included in
this subsystem.

Let $U$ be a finite propositional term generated by
\[
 U ::= a\mid b\mid c\mid F\mid\top\mid(U\to U).
\]
Replace the full four-input schema $L(a,b,c,u)$ by all instances of
$L(a,b,c,U(a,b,c,F,\top))$.  There is no bound on the size or depth of
$U$.

\paragraph{Theorem.}
For every such $U$, the resulting logical subsystem fails to derive at
least one instance of Peirce's law
\[
 ((P\to Q)\to P)\to P.
\]
In particular no implication-only reconstruction of the fourth
parameter, however large and whether or not it uses falsity, recovers
the original complete classical logical subsystem under these rules.

\subsection{One four-valued matrix covers every compound reconstruction}

Use values $\{0,1,2,3\}$, the sole designated value $3$, falsity $F=0$,
and the implication table
\[
\begin{array}{c|cccc}
 \to&0&1&2&3\\\hline
 0&3&3&3&3\\
 1&3&3&0&3\\
 2&3&3&3&3\\
 3&0&0&0&3
\end{array}.
\tag{M4}\label{eq:fourth-m4}
\]
Take a two-element individual domain.  Interpret equality by $3$ for
identical individuals and $0$ otherwise.  Interpret $\forall x$ by
minimum over the two assignments to $x$ in the order
\[
 0<2<1<3.
\]
This is a logical-matrix interpretation used for a syntactic
nonderivability proof; it is not being presented as ordinary Boolean
semantics for classical first-order logic.

\paragraph{Validity of the retained subsystem.}
The table respects modus ponens: if $A$ and $A\to B$ are designated,
then $A=3$ and row $3$ forces $B=3$.  Generalization preserves validity
because the minimum of two $3$'s is $3$.  The diagonal entries validate
identity and fresh atomic vacuity.  The all-$3$ zero row gives ex falso;
direct evaluation validates $K$ and double-negation elimination.

The universal-minimum interpretation satisfies quantifier distribution
for all $4^4$ pairs of two-point truth vectors.  Quantifier commutation
follows from commutation of finite minima.  If a formula has two-point
vector $(p_0,p_1)$ and $m=\min(p_0,p_1)$, then
$\exists x\forall x A$ has value $\neg\neg m$.  Direct evaluation
gives
\[
 (\neg\neg m)\to p_0=(\neg\neg m)\to p_1=3,
\]
which is the precise retained $\rho_6$, not a replacement quantifier
axiom.  Equality existence evaluates to $3$ for each possible equality
vector $(3,0),(0,3),(3,3)$.  The three equality schemas hold for all
assignments and all membership truth tables: a false equality
antecedent uses the all-$3$ zero row, and a true one identifies the
relevant entries.  Zero-diagonal membership validates irreflexivity.
The supplied checker independently exhausts these finite obligations.

\paragraph{The range argument.}
Evaluation of all $4^4$ tuples gives exactly three failures of the full
schema:
\[
 (a,b,c,u)=(0,3,1,2),\ (1,3,1,2),\ (2,3,1,2),
 \qquad L(a,b,c,u)=0.
\]
Every failure requires $u=2$.  But the range of the entire implication
table is $\{0,3\}$.  Consequently, if $U=V\to W$ is any compound term,
its value belongs to $\{0,3\}$, regardless of the values or sizes of
$V$ and $W$.  This proves validity of every compound specialization,
uniformly and without enumerating terms.  Equivalently, structural
induction on terms separates the five leaves from this single
implication-root case.

The four bare replacements $U=b,c,F,\top$ are also valid in
\eqref{eq:fourth-m4}.  For a failure one would need respectively
$b=2$, $c=2$, $0=2$, or $3=2$, while all displayed failures have
$b=3,c=1$.

Finally use the irreflexive membership table
\[
\begin{array}{c|cc}
 \in&0&1\\\hline
 0&0&1\\
 1&2&0
\end{array}.
\]
At the individual assignment $(v_0,v_1)=(0,1)$, let
$P=(v_0\in v_1)$ and $Q=(v_1\in v_0)$.  Then $P=1$, $Q=2$, and
\[
 P\to Q=0,\qquad (P\to Q)\to P=3,\qquad
 ((P\to Q)\to P)\to P=0.
\]
Every axiom is valid under every assignment and both inference rules
preserve validity, so this Peirce instance is not derivable.

\subsection{The one remaining leaf}

For the bare term $U=a$, use the three-valued matrix already present
in the preceding companion:
\[
\begin{array}{c|ccc}
 \to&0&1&2\\\hline
 0&2&2&2\\
 1&1&2&2\\
 2&0&1&2
\end{array},
\quad F=0,\quad\text{designated value }2.
\tag{M3}\label{eq:fourth-m3}
\]
Use the same two-element individual domain, crisp equality $2/0$,
and universal minimum in the order $0<1<2$.  All retained logical
schemas and the four additional schemas remain valid, as the separate
checker confirms.  The only failure of the full four-input $L$ is
\[
 (a,b,c,u)=(1,2,1,0),\qquad L=1.
\]
It cannot occur when $u=a$, so $L(a,b,c,a)$ is valid.
Give membership value $1$ at $(0,1)$ and $0$ at all other pairs.
At the assignment above take $P=(v_0\in v_1)$ and $Q=F$.
Peirce then has value $1$, which is undesignated.  This proves
nonderivability for the last leaf and completes the theorem.
\hfill$\square$

The numerical aliases of the source do not bypass the displayed
countermodels.  Code $1$ is interpreted as $v_0\in v_0=0$ under every
assignment, including assignments changed by quantifiers.  Code $0$
is truth, and the alias $\operatorname{Imp}(0,0)=0$ agrees with the
designated diagonal value $\top\to\top=\top$ in both matrices.

\subsection{Certificate, scope, and implications for compression}

Run the standard-library checker \file{check_fourth_input_no_go.py}.

Its output is \file{fourth_input_no_go_certificate.json}.

It contains the full tables, designated values, quantifier orders,
membership witnesses, failing tuples, and exact case counts.  It checks
the three equality rows against every membership table on the
two-element domain: $648$ assignments per row for M3 and $2048$ per
row for M4.  It separately checks all $81/256$ quantifier-distribution
tuples and all $9/16$ two-point instances of $\rho_6$.
This is exhaustive verification of finite algebraic facts.  The
unbounded term conclusion is the explicit range argument above,
not an extrapolation from a search-depth cutoff.

The three-valued matrix is inherited from the preceding companion's
positive-term obstruction.  The four-valued table and the partition
covering all terms are supplied in this report; no historical priority
is asserted for these matrices or the general use of matrix
countermodels in axiom research.

The theorem blocks a specific attractive optimization: replace the
fourth read by a propositional function of the other three inputs and
reuse the old completeness proof after deriving the missing classical
basis from the retained logical rows.  By itself, this unconditional
finite-matrix theorem does not rule out recovering the missing laws
from CP or GIC, changing the other logical schemas,
using a different interpretation of classical set theory, or employing
an alternative proof language.  CP and GIC were not checked in these
matrices.  Section~\ref{sec:colored-model-addendum} gives a conditional
extension to the full bundle for every covered replacement except bare $a$.
In particular neither inequiconsistency of the full modified
set theory nor a lower bound on binary Turing-machine states follows.

The finite class-theory, ground-rewriting, and intuitionistic routes in
the previous reports therefore remain separate compiler projects.
A fixed theory payload and a complete checker still have to be emitted
with their blank-tape initialization.  The construction in this
report takes the immediately measurable route of preserving the
already calibrated logical engine while reducing its binary control.
